\documentclass[../main.tex]{subfiles}
\ifSubfilesClassLoaded{\setcounter{chapter}{1}}{}
\begin{document}

\chapter{Jupyter, Kernel, Markdown; CRISP-DM Ringkas; Etika \& Lisensi Data}

\begin{subcpmk}
  \item \textbf{Sub-CPMK-P1:} Mengoperasikan \jupyter{}, mengelola kernel, menjelaskan etika dan lisensi penggunaan data pada tugas praktikum (selaras RPS Mg.\ 1).
\end{subcpmk}

\SesiRPS{1}{P1}{$\sim$170 menit}{CPMK-P1}

\section{Kemampuan akhir sesi}

Setelah sesi, mahasiswa mampu: (1) membuat notebook baru dan mengeksekusi sel Markdown serta sel kode Python; (2) menjelaskan enam fase CRISP-DM secara ringkas dan mengaitkannya dengan satu studi kasus; (3) mendokumentasikan sumber data, lisensi, dan batasan penggunaan di Markdown; (4) menyimpan berkas \texttt{.ipynb} sesuai konvensi penamaan.

\section{Teori dan bahan kajian}

\subsection{Sel, kernel, dan urutan eksekusi}

\begin{konsep}
\jupyter{} menggabungkan narasi (\textbf{sel Markdown}) dan kode (\textbf{sel kode}). \textbf{Kernel} adalah proses Python yang menjalankan kode. Urutan angka di sisi sel menunjukkan urutan eksekusi terakhir---bukan urutan di halaman. Setelah mengubah banyak sel, gunakan \textbf{Restart Kernel and Run All} untuk memverifikasi reproduktifitas.
\end{konsep}

\subsection{CRISP-DM dalam satu halaman}

\textit{Cross-Industry Standard Process for Data Mining} memetakan proyek data ke enam fase: \textbf{Business Understanding} (tujuan bisnis), \textbf{Data Understanding} (koleksi, kualitas), \textbf{Data Preparation} (cleaning, fitur), \textbf{Modeling}, \textbf{Evaluation}, \textbf{Deployment}. Pada praktikum 1 SKS, fokus kita adalah memahami bahwa notebook yang baik mencerminkan alur ini meskipun tidak semua fase dieksekusi penuh setiap minggu.

\subsection{Etika dan lisensi data}

\begin{catatan}
Cantumkan \textbf{URL sumber}, \textbf{lisensi} (CC, domain publik, syarat Kaggle, dll.), dan hindari mengunggah data yang berisi \textbf{PII} (data pribadi) tanpa persetujuan. Jika memakai data sintetis, tulis secara eksplisit bahwa data sintetis.
\end{catatan}

\section{Keselamatan kerja, etika, dan K3}

Tata tertib laboratorium; tidak makan/minum di dekat keyboard; matikan perangkat dengan benar; laporkan kabel longgar atau perangkat panas. Etika akademik: kutip tutorial dengan atribusi; tidak plagiasi notebook rekan.

\section{Sarana dan prasarana}

Python 3.10 atau lebih baru; \texttt{jupyter}, \texttt{notebook} atau JupyterLab; alternatif Google Colab; browser terkini; ruang penyimpanan untuk mengunduh \texttt{.ipynb}.

\section{Prosedur kerja (di lab)}

\begin{enumerate}[leftmargin=*]
  \item Verifikasi Python: buka terminal, jalankan perintah versi interpreter.
  \item Buka Jupyter atau VS Code dengan ekstensi Jupyter / unggah notebook ke Colab.
  \item Buat notebook baru: judul, NIM, nama, tanggal di sel Markdown pertama.
  \item Tulis paragraf CRISP-DM: untuk skenario yang Anda pilih, jelaskan minimal tiga fase yang paling relevan.
  \item Tambahkan sel kode yang mencetak string identitas dan operasi aritmetika sederhana.
  \item Cantumkan satu contoh tautan dataset terbuka beserta kalimat lisensi/atribusi.
  \item Simpan sebagai \path{Minggu01_NIM_Nama.ipynb} dan salin isi inti ke \path{notebooks/minggu_01.ipynb}.
\end{enumerate}

\section{Contoh kode dan studi kasus}

\subsection{Memeriksa versi Python di notebook}

\begin{lstlisting}[caption=Versi Python dan path executable]
import sys
print(sys.version)
print(sys.executable)
\end{lstlisting}

\subsection{Markdown vs kode}

\begin{lstlisting}[caption=Operasi sederhana di sel kode]
# Judul sel: perkenalan
nama_kampus = "Universitas Bale Bandung"
semester = 2026
print(f"Praktikum Data Science | {nama_kampus} | {semester}")
\end{lstlisting}

\subsection{Struktur dictionary untuk CRISP-DM ringkas}

\begin{lstlisting}[caption=Mapping fase CRISP-DM ke catatan proyek mini]
crisp_dm = {
    "business_understanding": "Mengurangi churn pelanggan e-commerce",
    "data_understanding": "Dataset transaksi terbuka, cek missing",
    "data_prep": "Imputasi, encoding kategori",
    "modeling": "Klasifikasi churn dengan RandomForest",
    "evaluation": "F1-score, confusion matrix",
    "deployment": "Dashboard internal (di luar cakupan praktikum)",
}
for fase, catatan in crisp_dm.items():
    print(f"{fase}: {catatan}")
\end{lstlisting}

\subsection{Menyimpan metadata sumber data (etika)}

\begin{lstlisting}[caption=Blok dokumentasi sumber data di notebook]
Sumber_data = {
    "url": "https://www.kaggle.com/datasets/...",
    "lisensi": "CC0 / sesuai halaman dataset",
    "tanggal_unduh": "2026-04-10",
    "catatan": "Tidak untuk komersial tanpa izin pemilik dataset.",
}
import json
print(json.dumps(Sumber_data, indent=2, ensure_ascii=False))
\end{lstlisting}

\section{Referensi sesi}

\begin{itemize}[leftmargin=*]
  \item Microsoft Learn --- Jupyter Notebook untuk Python. \url{https://learn.microsoft.com/en-us/training/modules/python-create-run-jupyter-notebook/}
  \item IBM --- CRISP-DM overview. \url{https://www.ibm.com/docs/en/spss-modeler/saas?topic=dm-crisp-help-overview}
  \item Kaggle Learn Python (landasan sintaks). \url{https://www.kaggle.com/learn/python}
  \item \texttt{sumber\_materi.md} di repositori modul.
\end{itemize}

\begin{aktivitas}
  \item Buat notebook dengan minimal \textbf{empat sel Markdown} dan \textbf{empat sel kode} yang saling mendukung narasi CRISP-DM.
  \item Jalankan \textbf{Restart Kernel and Run All}; perbaiki error hingga bersih.
  \item Tulis refleksi etika (150--250 kata) tentang risiko salahgunakan data tabular.
\end{aktivitas}

\begin{latihan}
  \item Apa beda kernel mati dan kernel sibuk? Bagaimana gejalanya di Jupyter?
  \item Sebutkan dua risiko jika sel dijalankan tidak berurutan saat demo ke dosen.
  \item Jelaskan kapan fase \textit{Evaluation} CRISP-DM paling erat hubungannya dengan metrik praktikum minggu 11--12.
\end{latihan}

\begin{asesmen}
\textbf{Formatif (observasi):} minimal tiga sel kode dan tiga sel Markdown bebas error; ada blok teks yang menyebut sumber data dan lisensi; \texttt{minggu\_01.ipynb} dapat di-\textit{run all}.

\textbf{Rubrik jejak:} kejelasan CRISP-DM; kelengkapan etika data; konsistensi nama berkas pengumpulan.
\end{asesmen}

\begin{checklist}
  \item Saya dapat membuat dan menjalankan notebook dari awal hingga akhir
  \item Saya menjelaskan CRISP-DM dengan bahasa saya sendiri
  \item Saya mendokumentasikan sumber dan lisensi data
\end{checklist}

\begin{rangkuman}
Minggu 1 menegaskan fondasi lingkungan Jupyter, dokumentasi CRISP-DM, dan etika data sebelum masuk ke Pandas.

\textbf{Kata kunci:} kernel, Markdown, CRISP-DM, lisensi, reproduktifitas
\end{rangkuman}

\end{document}
